\documentclass[10pt,a4paper]{amsart}
\usepackage[margin=19mm]{geometry}
\usepackage{amsmath,amssymb,mathtools}
\usepackage[scr=rsfs,cal=euler]{mathalfa}
\usepackage{xfrac}
\usepackage[unicode,hidelinks,bookmarks=false]{hyperref}
\newcommand{\ie}{i.e.\ }
\newcommand{\cf}{cf.\ }
\newcommand{\resp}{respectively\ }
\newcommand{\Subsection}{\subsection}
\providecommand{\nobreakdash}{\nobreak\hbox{-}}
\setlength{\emergencystretch}{2em}
\multlinegap0pt
\usepackage{bm}
\usepackage{bbm}
\usepackage{dsfont}
\usepackage{xspace}
\usepackage{cancel}
\usepackage{mathtools}
\usepackage{amsthm}
\makeatletter
\@ifundefined{theorem}{\newtheorem{theorem}{Theorem}}{}
\makeatother
\DeclareMathOperator{\rank}{rank}
\newcommand{\bbX}{\mathbf{X}}
\newcommand{\bphi}{\boldsymbol{\varphi}}
\newcommand{\trp}{^{\mathsf{T}}}  % adjoint (maybe use * or H instead?)
\title[eGAD! double descent is explained by Generalized Aliasing De]{eGAD! double descent is explained by Generalized Aliasing Decomposition}
\author{Mark K. Transtrum, Gus L. W. Hart, Tyler J. Jarvis, Jared P. Whitehead}
\date{Source: arXiv:2408.08294v4 (CC BY 4.0)}
\begin{document}
\maketitle

\noindent\emph{Source and licence.} eGAD! double descent is explained by Generalized Aliasing Decomposition. Version arXiv:2408.08294v4 (CC BY 4.0), distributed under the Creative Commons Attribution 4.0 International licence, as recorded in the arXiv licence metadata for this exact version and shown on the abstract page https://arxiv.org/abs/2408.08294v4. Full rights notice: licenses/arxiv-2408.08294-CC-BY-4.0.txt.\par\smallskip

\noindent\textit{Editorial scope.} Counted unit: the Theorem whose source label is \texttt{thm:add-column-inverse-norm} in the article (source0.tex lines 604--611, source label \texttt{thm:add-column-inverse-norm}), delivered together with the article's own complete original proof of it (source0.tex lines 612--640, one \texttt{proof} environment of 29 lines). No mathematics is written, repaired or re-derived: statement and proof are the article's own text line for line. Required context retained uncounted: thm:interlaced-eigenvalues (lines 591--596). Every definition, display and label of those lines is retained unchanged. Omitted from this package and not redistributed: 1--590, 597--603, 641--1197. Nothing inside a retained excerpt is omitted or altered: every retained body is delivered exactly as the article's lines are written, so no omission receipt of any kind accompanies this package. Citations: the retained excerpts carry no citation command at all, so no bibliography entry of the article is transcribed and no reference is redistributed. Reference adaptation: none was needed and none was made; every reference inside the delivered unit resolves inside it, so no unresolved reference or citation is produced. Printed slips kept literal: no printed word, symbol, hypothesis or number of the counted statement or of its proof was altered. Layout and class: the article's own class is \texttt{revtex4-2} with packages including fontenc, amsmath, bm, graphicx, amssymb, amsthm, hyperref, xcolor, sidecap, mathrsfs; the wrapper is the lane's pinned \texttt{amsart} editorial wrapper at 10pt on A4, which restates the theorem-like environments the retained text needs and adds the article's own macro definitions that the retained text needs (\texttt{\textbackslash bbX}, \texttt{\textbackslash bphi}, \texttt{\textbackslash rank}, \texttt{\textbackslash trp}), with extra wrapper packages (bm, bbm, dsfont, xspace, cancel, mathtools). The delivered numbering is the unit's own. Assets: no retained span contains a figure environment, a graphics-inclusion command, a PDF-page inclusion, a table or a TikZ picture; no image file is copied into this package, the package declares no asset dependency and no image file is redistributed with it. Licence: version arXiv:2408.08294v4 of this article is distributed under the Creative Commons Attribution 4.0 International licence, as recorded in the arXiv licence metadata for this exact version and shown on the abstract page https://arxiv.org/abs/2408.08294v4. A full rights notice is delivered at licenses/arxiv-2408.08294-CC-BY-4.0.txt. Characters: every retained excerpt is pure ASCII, so the delivered engine, XeLaTeX with its default Latin Modern text font, reports no missing character.
\begin{theorem}\label{thm:interlaced-eigenvalues}
Let $\Phi$ be an $n\times n$ Hermitian matrix with eigenvalues $\alpha_1\ge \alpha_2 \ge \cdots \ge \alpha_n$ and let $C$ be a positive semidefinite matrix of rank $1$.  The eigenvalues  $\beta_1\ge \beta_2 \ge \cdots \ge \beta_n$ of the matrix $\Xi = \Phi + C$ satisfy 
\[
\beta_1 \ge \alpha_1 \ge \beta_2 \ge \alpha_2 \ge \cdots \ge \beta_n \ge \alpha_n.
\]
\end{theorem}

\begin{theorem}\label{thm:add-column-inverse-norm}
Let $\bbX$ be an $n\times m$ matrix of rank $r$ with smallest  singluar value $\sigma_r>0$.  Let $\widetilde{\bbX} = [\bbX | \bphi]$ be the $n\times (m+1)$ matrix obtained by adjoining an $n$-dimensional column vector $\bphi$ to $\bbX$.  The smallest singular value $\widetilde{\sigma}_{\mathrm{min}}>0$ of $\widetilde{\bbX}$ satisfies the following relations:
\begin{align*}
0 & < \widetilde{\sigma}_{\mathrm{min}} \le \sigma_r  & 
\text{if $\rank(\bbX) < \rank(\widetilde{\bbX})$,}\\
0 & < {\sigma}_r \le \widetilde{\sigma}_{\mathrm{min}}  & \text{if $\rank(\bbX) = \rank(\widetilde{\bbX})$.}
\end{align*}
\end{theorem}

\begin{proof}
Both $\bbX \bbX\trp$ and $\widetilde{\bbX}\widetilde{\bbX}\trp$ are $n \times n$ positive definite Hermitian matrices.
The singular value decomposition of $\bbX$ shows that the singular values $\sigma_1\ge \cdots \ge \sigma_r >0$ and the eigenvalues $\lambda_1\ge \cdots \ge \lambda_m$ of $\bbX \bbX\trp$ satisfy 
\[
\lambda_1 =\sigma^2_1 \ge \lambda_2 = \sigma^2_1 \ge \cdots \ge \lambda_r = \sigma^2_r > 0 = \lambda_{r+1}.  
\]
Similarly, the singular values $\widetilde{\sigma}_1 \ge \widetilde{\sigma}_2 \ge \cdots$ and eigenvalues $\widetilde{\lambda}_1 \ge \cdots \ge \widetilde{\lambda}_m$ of $\widetilde{\bbX} \widetilde{\bbX}\trp$ satisfy
\[
\widetilde{\lambda}_1 = \widetilde{\sigma}^2_1 \ge \widetilde{\lambda}_2 = \widetilde{\sigma}^2_2 \cdots \ge \widetilde{\lambda}_r = \widetilde{\sigma}^2_r \ge  \widetilde{\lambda}_{r+1} \ge \cdots,
\]
where $\widetilde{\lambda}_{r+1} = 0$ if $\rank(\widetilde{\bbX}) = r$, but $\widetilde{\lambda}_{r+1} > 0$  if $\rank(\widetilde{\bbX}) = r+1$.
Expanding $\widetilde{\bbX}\widetilde{\bbX}\trp$ gives
$\widetilde{\bbX}\widetilde{\bbX}\trp = \bbX \bbX\trp + \bphi\bphi\trp$,
where $\bphi \bphi\trp$ is positive semidefinite, so Theorem~\ref{thm:interlaced-eigenvalues} implies that 
\[
\widetilde{\lambda}_1 \ge \lambda_1 \cdots \ge \lambda_{r-1} \ge \widetilde{\lambda}_r \ge \lambda_r \ge \widetilde{\lambda}_{r+1} \ge 0. 
\]

If $\rank(\widetilde{\bbX}) = r+1$ (that is, $\bphi$ is not in the column space of $\bbX$), then 
$
\lambda_r = \sigma^2_r \ge \widetilde{\lambda}_{r+1} = \widetilde{\sigma}^2_{r+1} >0$.
Taking square roots gives  
$\sigma_r > \widetilde{\sigma}_{r+1} = \widetilde{\sigma}_{\mathrm{min}} > 0$.

If $\rank(\widetilde{\bbX}) = r$ (that is, $\bphi$ is in the column space of $\bbX$), then the smallest nonzero eigenvalue of  $\widetilde{\bbX}\widetilde{\bbX}\trp$ is $\widetilde{\lambda}_r$, which satisfies  
$\lambda_{r-1} \ge \widetilde{\lambda}_r \ge \lambda_r > 0$. 
Taking square roots gives 
$\widetilde{\sigma}_{\mathrm{min}}=\widetilde{\sigma}_{r}  \ge \sigma_r > 0$, as required.
\end{proof}

\end{document}
